% SOURCE_AUTHORITY: sga5_fr_workpass.tex, lines 15246--15246 (LF-normalized unit).
% SOURCE_SHA256: 791F4EFFC5E02832D5D77ED03518C8156D6F07E4C8238B03545DB93D883FBB28
Para demonstrar (ii), começa-se observando que o caso em que $P_0$ é livre é evidente (sem hipótese alguma sobre o ideal $I$). Passa-se ao caso geral em que $P_0$ é projetivo introduzindo um $A_0$-módulo livre $L_0$ do qual $P_0$ é somando direto e um $A$-módulo livre $L$ tal que $T(L)\simeq L_0$. O módulo $P_0$ identifica-se com a imagem de um projetor $e_0:L_0\to L_0$, e é preciso ver que existe um projetor $e:L\to L$ que levanta $e_0$, o que resulta de Bourbaki, Alg. comm., cap.~III, §4, n\textsuperscript{o}~6, lema 2, aplicado à sub-$A$-álgebra $B$ de $\End(L)$ gerada por um endomorfismo de $L$ que levanta $e_0$ (cuja existência é garantida por (i)) e ao ideal (nilpotente) dado pela interseção de $B$ com o núcleo do homomorfismo canônico $\End(L)\to\End(L_0)$. Isto demonstra a existência de $P$ em (ii). Seja ainda $P'$ como no enunciado de (ii); então o isomorfismo composto $T(P)\simeq P_0\simeq T(P')$ provém, em virtude de (i), de um homomorfismo $f:P\to P'$. Como $P$ e $P'$ são planos sobre $A$ e $T(f)$ é um isomorfismo, conclui-se que $\operatorname{Gr}_I(f)$ é um isomorfismo, e portanto $f$ é um isomorfismo (pois $I$ é nilpotente), de onde se segue (ii).
